\documentclass[11pt]{article}

\usepackage[margin=1in]{geometry}
\usepackage{amsmath, amssymb}
\usepackage{graphicx}
\usepackage{booktabs}
\usepackage{caption}
\usepackage{hyperref}

\title{Replication of Gourio \& Miao (2010): \\
``Firm Heterogeneity and the Long-Run Effects of Dividend Tax Reform''}
\author{Robert Kirkby}
\date{\today}

\begin{document}
\maketitle

\section{Introduction}
This document replicates the main quantitative results of Gourio and Miao
(2010, \emph{AEJ:~Macro}). The model is a Bewley-style stationary equilibrium
of heterogeneous firms with idiosyncratic productivity shocks, capital
adjustment costs, corporate and personal taxation, and a constraint that
firms cannot repurchase shares. The exercise is to compute the long-run
effects of cutting the dividend tax (and the joint cut in dividend and
capital-gains taxes) on aggregate capital, output, consumption, and welfare.

All tables are generated by \texttt{GourioMiao2010.m} and live under
\texttt{./SavedOutput/LatexInputs/}; figures reproduced from the original
article are placed under \texttt{./SavedOutput/Graphs/} as image files.

\section{Model}
A representative household owns equity in a continuum of firms indexed by
idiosyncratic productivity $z$. Each firm operates a Cobb--Douglas technology
$y = z k^{\alpha_k} \ell^{\alpha_l}$, hires labor in a spot market at wage
$w$, and chooses next-period capital $k'$. Capital adjustment costs are
quadratic in net investment:
\[
\Phi(k',k) \;=\; \frac{\phi^{c}}{2}\,\frac{(k'-(1-\delta)k)^{2}}{k}.
\]
Corporate income tax $\tau_{\text{corp}}$ is levied on profits net of
depreciation expensing and a fraction $\phi$ of adjustment costs:
\[
T \;=\; \pi - \delta k - \phi\,\Phi(k',k), \qquad \pi = y - w\ell.
\]
After-tax cash flow is paid out as dividends $d \ge 0$ or by issuing new
equity $s \ge 0$; share repurchases are not allowed. Personal income, dividend,
and capital-gains taxes are $\tau_{i}$, $\tau_{d}$, and $\tau_{cg}$ respectively.
The household has preferences $U(C,L) = \log C - \tfrac{h}{2}L^{2}$, which
delivers the labor-market clearing condition $h C L = (1-\tau_{i}) w$.
A no-arbitrage condition on equity returns pins the firm discount factor:
\[
\hat{\beta} \;=\; \frac{1}{\,1 + r\,(1-\tau_{i})/(1-\tau_{cg})\,},
\qquad \beta(r(1-\tau_{i}) + 1) = 1.
\]

\paragraph{Firm problem (as written in GM2010).}
With $d$ and $k'$ both as decisions, the firm chooses $(d,k')$ to maximize
\[
F \;=\; \frac{1-\tau_{d}}{1-\tau_{cg}}\,d \;-\; s,
\qquad
s \;=\; d + I + \Phi(k',k) - \big(\pi - \tau_{\text{corp}} T\big),
\qquad I \equiv k' - (1-\delta)k,
\]
subject to $d \ge 0$ and $s \ge 0$, and the Bellman equation
\(
V(k,z) = \max_{d,k'} \{ F + \hat{\beta}\,\mathbb{E}[V(k',z') \mid z]\}.
\)

\section{Algebraic simplifications used in the code}
Two static substitutions reduce the firm problem to a one-control Bellman
equation in $k$. Both substitutions are exact and hold for any wage $w$ (in
or out of general equilibrium); they merely use the firm's interim FOCs
analytically rather than numerically.

\subsection{\texorpdfstring{Substitute out labor: $\pi = (1-\alpha_l)\,y$}{Substitute out labor}}
The static labor FOC gives $w = \alpha_l z k^{\alpha_k} \ell^{\alpha_l-1}$
and therefore $w\ell = \alpha_l y$. Profits are then
\[
\pi \;=\; y - w\ell \;=\; (1-\alpha_l)\,y.
\]
Solving for the optimal $\ell$ and substituting back yields closed forms for
output, labor demand, and profit in $(k, z; w)$:
\[
\ell^{\star}(k,z;w) \;=\; \Big(\tfrac{\alpha_l z k^{\alpha_k}}{w}\Big)^{\!1/(1-\alpha_l)},
\quad
y(k,z;w) \;=\; (z k^{\alpha_k})^{1/(1-\alpha_l)}\,(\alpha_l/w)^{\alpha_l/(1-\alpha_l)},
\quad
\pi(k,z;w) \;=\; (1-\alpha_l)\,y(k,z;w).
\]
This identity follows from \emph{firm} optimization at the labor stage; it
does not require labor-market clearing, so it is valid both in and out of GE.

\subsection{Eliminate the dividend as a decision}
Use the budget constraint $s = d - A$ with
\[
A(k',k,z;w,\boldsymbol{\theta})
\;\equiv\; \pi - \tau_{\text{corp}} T - I - \Phi(k',k)
\]
(the dividend the firm would pay if it issued no equity). Substituting into $F$:
\[
F \;=\; s\,\Big[\tfrac{1-\tau_{d}}{1-\tau_{cg}} - 1\Big] \;+\; \tfrac{1-\tau_{d}}{1-\tau_{cg}}\,A.
\]
With $\tau_{d} \ge \tau_{cg}$ (baseline: $\tau_{d}=0.25$, $\tau_{cg}=0.20$),
the coefficient on $s$ is non-positive, so the firm sets $s$ as small as
feasible. Combined with $d \ge 0$:
\begin{itemize}
  \item If $A \ge 0$: optimum is $d^{\star} = A,\ s^{\star} = 0$ and
        $F = \tfrac{1-\tau_{d}}{1-\tau_{cg}}\,A$.
  \item If $A < 0$: optimum is $d^{\star} = 0,\ s^{\star} = -A > 0$ and $F = A$.
\end{itemize}
In one line,
\[
\boxed{\;F(k',k,z;w,\boldsymbol{\theta})
\;=\; \tfrac{1-\tau_{d}}{1-\tau_{cg}}\,\max\{A,0\} \;+\; \min\{A,0\},\;}
\]
with $d^{\star}(k',k,z) = \max\{A,0\}$ and $s^{\star}(k',k,z) = \max\{-A,0\}$.
The dividend is therefore not an independent control: given any policy for
$k'$, the dividend and equity-issuance policies are pinned down algebraically.

\subsection{The simplified Bellman equation}
After both substitutions, the firm's recursive problem is a single Bellman
equation in $k$ with the single control $k'$:
\[
V(k,z) \;=\; \max_{k' \in \mathcal{K}}\left\{
\tfrac{1-\tau_{d}}{1-\tau_{cg}}\,\max\{A,0\} + \min\{A,0\}
\;+\; \hat{\beta}\,\mathbb{E}\big[V(k',z')\,\big|\,z\big]\right\},
\]
where
\[
A(k',k,z;w,\boldsymbol{\theta})
\;=\; (1-\tau_{\text{corp}})(1-\alpha_l)\,y(k,z;w)
\;+\; \tau_{\text{corp}}\delta k
\;-\; I
\;-\; (1-\tau_{\text{corp}}\phi)\,\Phi(k',k).
\]
Numerically this is what \texttt{GourioMiao2010.m} solves: it sets
$n_{d}=0$, gives the toolkit a single decision grid (capital), and never
discretizes dividends. The optimal $d^{\star} = \max\{A,0\}$ and
$s^{\star} = \max\{-A,0\}$ are recovered from $k'$-policy at the post-processing
stage for table-building (see \texttt{GourioMiao2010\_DividendFn.m} and
\texttt{GourioMiao2010\_NewEquityFn.m}).

\paragraph{Financing regimes.}
Because exactly one of $d^{\star}$, $s^{\star}$ is non-zero, the three GM2010
financing regimes correspond cleanly to the sign of $A$:
\[
\underbrace{A > 0}_{\text{dividend regime}}
\qquad
\underbrace{A \approx 0}_{\text{liquidity-constrained}}
\qquad
\underbrace{A < 0}_{\text{equity-issuance regime}}.
\]
The liquidity-constrained mass concentrates at the kink in $F$ where the slope
in $k'$ changes; with a discrete $a$-grid the optimum sits on a grid point,
so the regime is identified with $|A| \le \text{tol}$ (we use the
$a$-grid spacing).

\section{Calibration}
The baseline parametrization (Figure~\ref{fig:GM2010_Table3}) matches the
values reported in GM2010.

\begin{figure}[!ht]\centering
\includegraphics[width=0.85\textwidth]{./OriginalMaterials/GM2010_OriginalTable3.png}\\[1em]
\input{./SavedOutput/LatexInputs/GourioMiao2010_Table3.tex}
\caption{Baseline parametrization. Top: original from Gourio \& Miao (2010). Bottom: replication.}\label{fig:GM2010_Table3}
\end{figure}

Productivity is discretized with a 10-point Tauchen--Hussey grid; capital
uses $n_a = 701$ evenly spaced points on $[0.001,\,8]$.

\section{Baseline results}
With the wage solved from the household labor-market FOC,
Figures~\ref{fig:GM2010_Table4} and~\ref{fig:GM2010_Table5} report the
aggregate and cross-sectional moments and the share of firms in each
financing regime, each alongside the corresponding original.

\begin{figure}[!ht]\centering
\includegraphics[width=0.85\textwidth]{./OriginalMaterials/GM2010_OriginalTable4.png}\\[1em]
\input{./SavedOutput/LatexInputs/GourioMiao2010_Table4.tex}
\caption{Aggregate and cross-sectional moments. Top: original from Gourio \& Miao (2010). Bottom: replication.}\label{fig:GM2010_Table4}
\end{figure}

\begin{figure}[!ht]\centering
\includegraphics[width=0.85\textwidth]{./OriginalMaterials/GM2010_OriginalTable5.png}\\[1em]
\input{./SavedOutput/LatexInputs/GourioMiao2010_Table5.tex}
\caption{Distribution of firms across finance regimes. Top: original from Gourio \& Miao (2010). Bottom: replication.}\label{fig:GM2010_Table5}
\end{figure}

\section{Dividend tax reform}
The reform exercise considers four pairs of $(\tau_{d}, \tau_{cg})$ relative
to the baseline. Figure~\ref{fig:GM2010_Table6} gives the long-run percentage
changes in aggregates; Figure~\ref{fig:GM2010_Table7} reports how the firm
distribution across financing regimes shifts under the small ($\tau_d=0.22$)
reform; Figure~\ref{fig:GM2010_Table8} decomposes the productivity gain.

\begin{figure}[!ht]\centering
\includegraphics[width=0.85\textwidth]{./OriginalMaterials/GM2010_OriginalTable6.png}\\[1em]
\input{./SavedOutput/LatexInputs/GourioMiao2010_Table6.tex}
\caption{Aggregate effects of the dividend tax reform. Top: original from Gourio \& Miao (2010). Bottom: replication.}\label{fig:GM2010_Table6}
\end{figure}

\begin{figure}[!ht]\centering
\includegraphics[width=0.85\textwidth]{./OriginalMaterials/GM2010_OriginalTable7.png}\\[1em]
\input{./SavedOutput/LatexInputs/GourioMiao2010_Table7.tex}
\caption{Distribution of firms across finance regimes after the small reform. Top: original from Gourio \& Miao (2010). Bottom: replication.}\label{fig:GM2010_Table7}
\end{figure}

\begin{figure}[!ht]\centering
\includegraphics[width=0.85\textwidth]{./OriginalMaterials/GM2010_OriginalTable8.png}\\[1em]
\input{./SavedOutput/LatexInputs/GourioMiao2010_Table8.tex}
\caption{Productivity gains from the dividend tax cut. Top: original from Gourio \& Miao (2010). Bottom: replication.}\label{fig:GM2010_Table8}
\end{figure}

\section{Sensitivity}
Figures~\ref{fig:GM2010_Table9} and~\ref{fig:GM2010_Table10} re-run the
baseline moments and the $(\tau_{d},\tau_{cg})=(0.15,0.15)$ reform for
alternative values of $\rho_{z}$, $\sigma_{z}$, and the adjustment-cost
parameter $\phi^{c}$.

\begin{figure}[!ht]\centering
\includegraphics[width=0.85\textwidth]{./OriginalMaterials/GM2010_OriginalTable9.png}\\[1em]
\input{./SavedOutput/LatexInputs/GourioMiao2010_Table9.tex}
\caption{Moments for different parameter values. Top: original from Gourio \& Miao (2010). Bottom: replication.}\label{fig:GM2010_Table9}
\end{figure}

\begin{figure}[!ht]\centering
\includegraphics[width=0.85\textwidth]{./OriginalMaterials/GM2010_OriginalTable10.png}\\[1em]
\input{./SavedOutput/LatexInputs/GourioMiao2010_Table10.tex}
\caption{Reform effects under different parameter values. Top: original from Gourio \& Miao (2010). Bottom: replication.}\label{fig:GM2010_Table10}
\end{figure}

\section{Notes on the code}
\begin{itemize}
\item \texttt{GourioMiao2010\_ReturnFn.m} implements the simplified return
function $F = \tfrac{1-\tau_{d}}{1-\tau_{cg}}\max\{A,0\} + \min\{A,0\}$, with
$y$ and $\pi$ in closed form. The driver sets $n_{d}=0$ and \(d\_grid=[]\).
\item \texttt{GourioMiao2010\_DividendFn.m} and
\texttt{GourioMiao2010\_NewEquityFn.m} return $d^{\star}=\max\{A,0\}$ and
$s^{\star}=\max\{-A,0\}$ respectively; both are used only at the
post-processing stage to build tables.
\item Regime classification uses $\text{tol}$ equal to the $a$-grid spacing,
applied to $d^{\star}$ and $s^{\star}$; no separate tolerance on dividends is
needed because $d^{\star}$ and $s^{\star}$ are computed analytically.
\end{itemize}

\end{document}
